\documentclass[11pt]{article}
\usepackage[margin=1in]{geometry}
\usepackage{amsmath,amssymb,amsthm}
\usepackage{enumitem}

\theoremstyle{plain}
\newtheorem{theorem}{Theorem}
\newtheorem{lemma}[theorem]{Lemma}
\theoremstyle{definition}
\newtheorem{definition}[theorem]{Definition}

\begin{document}
\pagestyle{empty}

\begin{definition}[The two rules used here]
\label{def:rules}
Every vertex of a graph $G$ carries a status in $\{\texttt{IN}, \texttt{OUT},
\texttt{UNKNOWN}\}$, initially \texttt{UNKNOWN}. Two rules are applied, repeatedly
and in any order, each firing at most once per vertex, until neither can fire:
\begin{itemize}[nosep]
\item \textbf{(Spread)} If $v$ is \texttt{IN}, mark every neighbour of $v$
\texttt{OUT} (immediately — not deferred to a later pass).
\item \textbf{(Force)} If $v$ is \texttt{UNKNOWN} and $v$'s neighbours not
currently marked \texttt{OUT} contain no independent pair (equivalently, that set
is a clique, possibly of size $\le 1$), mark $v$ \texttt{IN}.
\end{itemize}
A vertex left \texttt{UNKNOWN} at a fixed point is one neither rule ever applied to.
\end{definition}

\begin{lemma}[Soundness]
\label{lem:sound}
Whenever this process fixes a vertex $v$ as \texttt{IN} or \texttt{OUT}, that status
holds for \emph{every} 2-kernel of $G$ — not merely for whatever candidate the
process itself is building.
\end{lemma}

\begin{proof}
Strong induction on the order in which vertices get fixed. Suppose the claim holds
for every vertex fixed strictly before $v$, and consider how $v$ was fixed.

\emph{Case: $v$ fixed \texttt{OUT}.} This happens via Spread, triggered by some
neighbour $u$ already \texttt{IN}, with $u$ fixed strictly before $v$. By the
inductive hypothesis $u \in S$ for every 2-kernel $S$; since $S$ is independent and
$u$ is a neighbour of $v$, $v \notin S$ for every such $S$.

\emph{Case: $v$ fixed \texttt{IN}.} This happens via Force: at that moment, $v$'s
neighbours not marked \texttt{OUT} contain no independent pair. Every neighbour of
$v$ already marked \texttt{OUT} was fixed strictly before $v$, so by the inductive
hypothesis it is excluded from every 2-kernel $S$. Take any 2-kernel $S$ of $G$ and
suppose toward a contradiction $v \notin S$. Then $S \cap N(v)$ avoids every
neighbour already known \texttt{OUT}, so $S \cap N(v)$ sits entirely inside $v$'s
current non-\texttt{OUT} neighbours. Since $v \notin S$, 2-domination requires
$|S \cap N(v)| \ge 2$ — but $S$ is independent, so those $\ge 2$ vertices are
pairwise non-adjacent, i.e.\ an independent pair inside $v$'s non-\texttt{OUT}
neighbourhood. This contradicts the very condition that made Force fire. So
$v \in S$ for every 2-kernel $S$.
\end{proof}

\begin{theorem}[E2.1]
On a chordal graph $G$, running (Spread) and (Force) to a fixed point never leaves a
vertex \texttt{UNKNOWN}. Consequently $G$ has at most one 2-kernel, and 2-KERNEL is
solvable in polynomial time on chordal graphs (the general problem is NP-complete).
\end{theorem}

\begin{proof}
Run the process to a fixed point and suppose toward a contradiction that some vertex
remains \texttt{UNKNOWN}; let $W \ne \emptyset$ be the set of all such vertices.

No vertex of $W$ has an \texttt{IN} neighbour: if $v \in W$ had one, Spread would
already have fired on it the moment that neighbour became \texttt{IN}, marking $v$
\texttt{OUT} — contradicting $v \in W$. So for every $v \in W$, the neighbours of $v$
not marked \texttt{OUT} are exactly $N(v) \cap W$.

$G[W]$, the subgraph induced on $W$, is an induced subgraph of the chordal graph $G$,
hence itself chordal (a classical, standard fact about chordal graphs). Since
$W \ne \emptyset$, $G[W]$ has a simplicial vertex $x$ — by definition, $x$'s
neighbourhood \emph{within $G[W]$}, namely $N(x) \cap W$, is a clique, i.e.\ contains
no independent pair (existence of a simplicial vertex in every non-empty chordal
graph is likewise classical). By the previous paragraph, $N(x) \cap W$ is exactly
$x$'s current non-\texttt{OUT} neighbourhood — so Force applies to $x$, contradicting
the assumption that we are at a fixed point where no rule can fire.

Hence $W = \emptyset$: every vertex ends up \texttt{IN} or \texttt{OUT}. By
Lemma~\ref{lem:sound}, this final assignment records, for every vertex, its status in
\emph{every} 2-kernel of $G$. So if $G$ has a 2-kernel at all, it is exactly
$S = \{v : v \text{ fixed } \texttt{IN}\}$, and no other set can be one — $G$ has at
most one 2-kernel.

For the complexity claim: chordal graphs are recognized, and a simplicial-vertex
elimination ordering produced, in polynomial time (via lexicographic breadth-first
search); each application of Force or Spread is a polynomial-time clique/neighbour
check; and the process above fixes one previously-\texttt{UNKNOWN} vertex per
firing, so it terminates after at most $n$ firings. Checking the resulting $S$
against the definition of a 2-kernel is also polynomial. So the whole procedure —
decide existence, and produce the unique 2-kernel if one exists — runs in
polynomial time.
\end{proof}

\end{document}
